\section{Proof of the separation}\label{sec:separation}

We now combine the invariant query, support bound, counting transfer and
exact computation formulas. The constants controlling supports and formula
width are fixed by the program before the grid is chosen.

\begin{proof}[Proof of Theorem~\ref{thm:main}]
The fixed eight-symbol vocabulary and the query $Q$ are given in
Section~\ref{sec:query}. Proposition~\ref{prop:query-ptime} proves that $Q$
is isomorphism-invariant and belongs to deterministic polynomial time.
Suppose it were definable in full $\CPT$ with counting. The interpretation
characterization in Section~\ref{sec:coding} would give a fixed deciding
interpretation program $P$.

Choose its polynomial bound $B=B(N)$ and the constant $q$
from the trace bound, and choose the constant $m$ from the syntax.
Take $L$ sufficiently large that the trace bound $9B^3\le N^q$ and
the width estimates below hold, and put $M=\lfloor L^{7/4}\rfloor$.
Corollary~\ref{cor:transitive-support} places both encoded traces in
their supported domains with
\[
 s=O_q\bigl(L(1+\log L)^2\bigr).
\]
Proposition~\ref{prop:base-equivalence}
gives base $C^M$ equivalence, and
Proposition~\ref{prop:base-homogeneity} gives homogeneity for tuples of
length at most $\max\{2,m\}s$: its width bound is
$O_{q,m}(L^{3/2}(1+\log L)^3)=o(L^{7/4})$.
We also have $\max\{4,m+1\}s\le M$.
These are exactly the hypotheses of Theorem~\ref{thm:hf-transfer}, so the
two supported domains agree on all $C^m$ sentences.

Use the same $B$, ordinal indices, and syntactic definitions on both
sides. Up to the earlier halt stage both sets of state definitions are
exact. At that stage the common translated halt sentence is true on
one side, hence also on the other; thus both computations halt there.
Their common translated output sentence has the same truth value on
both sides. No deciding interpretation program can therefore give
opposite answers on the two structures.

This contradicts Proposition~\ref{prop:opposite-answers}, which gives
$Q(A_{b^0})$ and $\neg Q(A_{b^1})$ for every $L\ge2$. Therefore no such
program exists, proving the theorem.
\end{proof}
